% Book-local P12.2/P12.3 curve-level ideal-psh construction.
% The superscript ``id'' records that this is the closure of the displayed
% ideal-model family; no comparison with a separate nef or CLD psh cone is
% asserted.
\begin{theorem}[Ideal-psh completion and envelopes on reduced curve links]
\label{thm:p122-curve-ideal-psh-functoriality}
Let \((R,\mathfrak m)\) be a one-dimensional reduced Nagata noetherian
local ring, let \(\mathfrak a\subset R\) be \(\mathfrak m\)-primary, and
put \(V_{\mathfrak a}:=\mathcal V_{\mathfrak a}(R)\).  With the
principal-ideal model family
\(\mathscr P_{\mathfrak a}^{\mathrm{pr}}\) of
\eqref{eq:p122-principal-ideal-profile-family}, write
\(u^{\mathfrak a}_{f,N}\) when its normalizing ideal must be visible, and
define
\[
 \operatorname{CPSH}^{\mathrm{id}}_{\mathfrak a}(R)
 :=\overline{\mathscr P_{\mathfrak a}^{\mathrm{pr}}}
                         ^{\,\|\cdot\|_\infty},
 \qquad
 \operatorname{PSH}^{\mathrm{id},\bot}_{\mathfrak a}(R)
 :=\left\{\lim_{j\to\infty}u_j:
       u_j\in\operatorname{CPSH}^{\mathrm{id}}_{\mathfrak a}(R),\
       u_j\searrow\right\}.
 \tag{\(\triangle\)}\label{eq:p122-curve-ideal-psh-definition}
\]
The second class retains the identically \(-\infty\) bottom element;
removing it gives
\(\operatorname{PSH}^{\mathrm{id}}_{\mathfrak a}(R)\).  Then:
\begin{enumerate}
\item The two classes have the intrinsic descriptions
      \[
       \operatorname{CPSH}^{\mathrm{id}}_{\mathfrak a}(R)
          =\{u:V_{\mathfrak a}\to\mathbf R:u\leq0\},
       \qquad
       \operatorname{PSH}^{\mathrm{id},\bot}_{\mathfrak a}(R)
          =\{u:V_{\mathfrak a}\to[-\infty,0]\}.
       \tag{\(\square\)}\label{eq:p122-curve-ideal-psh-orthants}
      \]
      In particular, the continuous class is a closed convex cone stable
      under finite maxima, and the completed class is a complete lattice.
      Replacing the principal model family by all positive rational
      multiples of \(\log|\mathfrak b|\), with \(\mathfrak b\) ranging
      over the \(\mathfrak m\)-primary ideals, gives the same closure.
\item For an arbitrary extended-real obstacle
      \(\phi:V_{\mathfrak a}\to[-\infty,+\infty]\), its completed
      ideal-psh envelope
      \[
       P^{\mathrm{id}}_{\mathfrak a}\phi
       :=\sup\{u\in
          \operatorname{PSH}^{\mathrm{id},\bot}_{\mathfrak a}(R):
          u\leq\phi\}
      \]
      is attained and is given coordinatewise by
      \[
                       P^{\mathrm{id}}_{\mathfrak a}\phi
                              =\min\{\phi,0\}.
       \tag{\(\lozenge\)}\label{eq:p123-curve-ideal-psh-envelope}
      \]
      For finite \(\phi\), it belongs to
      \(\operatorname{CPSH}^{\mathrm{id}}_{\mathfrak a}(R)\), and the
      same function is the supremum of principal-ideal competitors in
      \eqref{eq:p123-principal-ideal-envelope}.
\item These constructions are independent of the chosen primary-ideal
      cross-section with an explicit homogeneous realization.  Put
      \[
       \mathcal C_R^{\mathrm{br}}
          :=\{t w_{\mathfrak q}:t>0,\
                  \mathfrak q\in\operatorname{Br}(R,\mathfrak m)\}.
      \]
      Restriction to \(V_{\mathfrak a}\) identifies the completed class
      in \eqref{eq:p122-curve-ideal-psh-orthants} with the nonpositive
      degree-one homogeneous functions
      \[
       G:\mathcal C_R^{\mathrm{br}}\longrightarrow[-\infty,0],
       \qquad G(t\nu)=tG(\nu).
       \tag{\(\circ\)}\label{eq:p122-curve-homogeneous-branch-psh}
      \]
      Finite-valued homogeneous functions correspond to the continuous
      class.  In particular, this homogeneous object is independent of
      \(\mathfrak a\).  For another
      \(\mathfrak m\)-primary ideal \(\mathfrak b\), let
      \(H_{\mathfrak a,\mathfrak b}\) be
      \eqref{eq:p121-change-normalizing-ideal} and define
      \[
       (T_{\mathfrak a,\mathfrak b}u)
          \bigl(H_{\mathfrak a,\mathfrak b}(v)\bigr)
          :=\frac{u(v)}{v(\mathfrak b)}.
       \tag{\(\diamondsuit\)}\label{eq:p122-curve-psh-normalization-transport}
      \]
      Division by a positive number fixes both infinities.  The map
      \(T_{\mathfrak a,\mathfrak b}\) is an order and lattice
      isomorphism on both classes in
      \eqref{eq:p122-curve-ideal-psh-definition}, carries each
      principal-ideal profile to the profile of the same pair \((f,N)\),
      extends by the same formula to extended-real obstacles, and satisfies
      \[
       T_{\mathfrak b,\mathfrak c}T_{\mathfrak a,\mathfrak b}
          =T_{\mathfrak a,\mathfrak c},
       \qquad
       T_{\mathfrak a,\mathfrak b}
          (P^{\mathrm{id}}_{\mathfrak a}\phi)
          =P^{\mathrm{id}}_{\mathfrak b}
             (T_{\mathfrak a,\mathfrak b}\phi).
       \tag{\(\clubsuit\)}\label{eq:p122-curve-psh-transport-cocycle}
      \]
\item Let
      \(\varphi:(R,\mathfrak m)\to(S,\mathfrak n)\) be a finite local
      homomorphism between one-dimensional reduced Nagata noetherian local
      rings, and let
      \(\rho_{\mathfrak a}:V_{\mathfrak aS}\to V_{\mathfrak a}\) be the
      restriction map of \Cref{prop:p126-finite-reduced-curve-link}.  Then
      \(\rho_{\mathfrak a}^*u=u\circ\rho_{\mathfrak a}\) preserves the
      continuous and completed ideal-psh classes (including the bottom
      element), sends
      \[
                u^{\mathfrak a}_{f,N}
                    \longmapsto
                u^{\mathfrak aS}_{\varphi(f),N},
      \]
      and commutes with completed envelopes:
      \[
       \rho_{\mathfrak a}^*
          (P^{\mathrm{id}}_{\mathfrak a}\phi)
        =P^{\mathrm{id}}_{\mathfrak aS}
          (\rho_{\mathfrak a}^*\phi).
       \tag{\(\spadesuit\)}\label{eq:p123-curve-psh-finite-map-envelope}
      \]
      Pullback is compatible with composition and with the transports
      \eqref{eq:p122-curve-psh-normalization-transport}; explicitly, for
      every second \(\mathfrak m\)-primary ideal \(\mathfrak b\),
      \[
       \rho_{\mathfrak b}^*T^R_{\mathfrak a,\mathfrak b}
          =T^S_{\mathfrak aS,\mathfrak bS}
             \rho_{\mathfrak a}^*.
      \tag{\(\bigstar\)}\label{eq:p122-curve-psh-map-normalization-square}
      \]
      On either class, the image of \(\rho_{\mathfrak a}^*\) consists
      exactly of the functions on \(V_{\mathfrak aS}\) that are constant on
      every fiber of \(\rho_{\mathfrak a}\).  Consequently,
      \(\rho_{\mathfrak a}^*\) is injective if and only if
      \(\rho_{\mathfrak a}\) is surjective, and it is surjective if and only
      if \(\rho_{\mathfrak a}\) is injective.  When
      \(\rho_{\mathfrak a}\) is surjective, pullback on the continuous class
      is an isometry:
      \[
       \|\rho_{\mathfrak a}^*u-\rho_{\mathfrak a}^*v\|_\infty
          =\|u-v\|_\infty.
      \]
      On the completed lattices, pullback has the explicit left and right
      adjoints
      \[
       (\rho_{\mathfrak a,!}w)(x)
          :=\sup\bigl(\{w(y):\rho_{\mathfrak a}(y)=x\}
                        \cup\{-\infty\}\bigr),
      \]
      \[
       (\rho_{\mathfrak a,*}w)(x)
          :=\inf\bigl(\{w(y):\rho_{\mathfrak a}(y)=x\}
                        \cup\{0\}\bigr),
      \]
      so that
      \[
       \rho_{\mathfrak a,!}w\leq u
          \Longleftrightarrow
       w\leq\rho_{\mathfrak a}^*u,
       \qquad
       \rho_{\mathfrak a}^*u\leq w
          \Longleftrightarrow
       u\leq\rho_{\mathfrak a,*}w.
       \tag{\(\heartsuit\)}\label{eq:p122-curve-psh-finite-map-adjoints}
      \]
      If \(\rho_{\mathfrak a}\) is surjective, both composites with
      pullback on the right are the identity on the base lattice; if it is
      injective, both composites with pullback on the left are the identity
      on the source lattice.  For arbitrary \(\rho_{\mathfrak a}\),
      \(\rho_{\mathfrak a}^*\rho_{\mathfrak a,*}w\) is the greatest
      fiber-constant minorant of \(w\), while
      \(\rho_{\mathfrak a}^*\rho_{\mathfrak a,!}w\) is its least
      fiber-constant majorant.  Thus their common fixed-point set is exactly
      the pullback image.
\end{enumerate}
Thus the ideal-closure candidate has a complete psh-style order, envelope,
normalization, and pullback calculus on reduced Nagata curve links.  The
superscript \({\rm id}\) is a scope marker: naming this closure does not by
itself identify it with a separately defined nef-b-divisor, CLD, or
complex-analytic psh class, and hence does not by itself retire P12.2 or
P12.3.  The theorem constructs no local intersection product or energy.
\end{theorem}

\begin{proof}
The link is finite and discrete by
\Cref{thm:p121-reduced-nagata-curve-link}.  The first equality in
\eqref{eq:p122-curve-ideal-psh-orthants} is exactly the uniform-closure
calculation in
\eqref{eq:p122-principal-ideal-orthant}.  If
\(u:V_{\mathfrak a}\to[-\infty,0]\), then
\[
                            u_j:=\max\{u,-j\}
\]
is finite, nonpositive, and decreases pointwise to \(u\).  This proves the
second equality.  The asserted cone and lattice properties are now
coordinatewise.  Every \(\mathfrak m\)-primary ideal function has strictly
negative rational coordinates, so
\eqref{eq:p122-principal-ideal-orthant} realizes the same function by a
principal model; this proves the last assertion of the first item.

The function \(w=\min\{\phi,0\}\) belongs to the completed class, lies
below \(\phi\), and dominates every nonpositive competitor.  This proves
\eqref{eq:p123-curve-ideal-psh-envelope}.  If \(\phi\) is finite, then
\(w\) is finite and hence continuous on the finite link.  Its equality with
the raw principal-ideal supremum is
\eqref{eq:p123-principal-ideal-envelope-formula}.

Write
\(v_{\mathfrak q}^{\mathfrak a}
  =w_{\mathfrak q}/d_{\mathfrak q}(\mathfrak a)\).
Every nonpositive homogeneous \(G\) restricts to a member of the completed
class by \eqref{eq:p122-curve-ideal-psh-orthants}.  Conversely, a function
\(u\) in that class has the unique homogeneous extension
\[
 G_u(tw_{\mathfrak q})
    :=t d_{\mathfrak q}(\mathfrak a)
          u(v_{\mathfrak q}^{\mathfrak a}).
\]
The same formula preserves finite values.  This proves
\eqref{eq:p122-curve-homogeneous-branch-psh}.  Moreover,
\[
 v_{\mathfrak q}^{\mathfrak a}(\mathfrak b)
   =\frac{d_{\mathfrak q}(\mathfrak b)}
          {d_{\mathfrak q}(\mathfrak a)},
 \qquad
 H_{\mathfrak a,\mathfrak b}
       (v_{\mathfrak q}^{\mathfrak a})
   =v_{\mathfrak q}^{\mathfrak b}.
\]
Hence \eqref{eq:p122-curve-psh-normalization-transport} is multiplication
by one positive scalar on each coordinate.  It is therefore an order and
lattice isomorphism, preserves decreasing limits, and has inverse
\(T_{\mathfrak b,\mathfrak a}\).  Moreover,
\[
 \frac{u^{\mathfrak a}_{f,N}(v)}{v(\mathfrak b)}
   =-\frac{v(f)}{Nv(\mathfrak b)}
   =u^{\mathfrak b}_{f,N}
       \bigl(H_{\mathfrak a,\mathfrak b}(v)\bigr).
\]
For a third ideal \(\mathfrak c\), putting
\(w=H_{\mathfrak a,\mathfrak b}(v)\) gives
\(w(\mathfrak c)=v(\mathfrak c)/v(\mathfrak b)\); the two denominators
cancel and prove the first identity in
\eqref{eq:p122-curve-psh-transport-cocycle}.  The second follows from
\(T(\min\{\phi,0\})=\min\{T\phi,0\}\).

The restriction map \(\rho_{\mathfrak a}\) is a map of finite sets, so
composition with it preserves finite nonpositive functions and pointwise
decreasing limits.  If \((f)\) is \(\mathfrak m\)-primary and principal,
then \(f\) avoids every minimal prime of \(R\).  The proof of
\Cref{prop:p126-finite-reduced-curve-link} shows that every minimal prime
of \(S\) contracts to a minimal prime of \(R\).  Thus \(\varphi(f)\)
avoids every minimal prime of the reduced ring \(S\), so it is a
nonzerodivisor.  It lies in \(\mathfrak n\), and in the one-dimensional
local ring \(S\) its principal ideal is therefore \(\mathfrak n\)-primary.
For \(w\in V_{\mathfrak aS}\),
\[
 (\rho_{\mathfrak a}^*u^{\mathfrak a}_{f,N})(w)
       =-\frac{(w\circ\varphi)(f)}N
       =-\frac{w(\varphi(f))}N
       =u^{\mathfrak aS}_{\varphi(f),N}(w),
\]
which proves the model formula.  Equation
\eqref{eq:p123-curve-psh-finite-map-envelope} follows by pulling back the
coordinatewise minimum formula.  Associativity of function composition
gives compatibility in a tower.  Finally,
\(w(\mathfrak bS)=(w\circ\varphi)(\mathfrak b)\) by
\eqref{eq:p126-extended-ideal-value}; inserting this equality into
\eqref{eq:p122-curve-psh-normalization-transport} proves
\eqref{eq:p122-curve-psh-map-normalization-square}.

The descriptions in \eqref{eq:p122-curve-ideal-psh-orthants} reduce the
remaining claims to finite sets.  A target function is a pullback precisely
when it is constant on every fiber: necessity is immediate, and for
sufficiency define the source function on the image of
\(\rho_{\mathfrak a}\) by the common fiber value and give it the value zero
at every missed point.  If \(\rho_{\mathfrak a}\) is surjective, equality of
two pullbacks implies equality of the source functions; if it is not
surjective, changing a source function only at a missed point disproves
injectivity.  If \(\rho_{\mathfrak a}\) is injective, the fiber condition is
vacant and every target function lifts; if it is not injective, a target
function taking the values zero and \(-1\) at two points in the same fiber
does not lift.  These witnesses are finite-valued, so the criteria hold for
both classes.  Finally, when \(\rho_{\mathfrak a}\) is surjective, taking
the maximum over its domain samples every point of its target, which proves
the displayed sup-norm identity.  For the completed classes, the two
fiberwise formulas take values in \([-\infty,0]\); the adjoined values
\(-\infty\) and zero are exactly the supremum and infimum over an empty
fiber.  The first equivalence in
\eqref{eq:p122-curve-psh-finite-map-adjoints} says coordinatewise that
every value of \(w\) over \(x\) is at most \(u(x)\), and the second says
that \(u(x)\) is at most every value of \(w\) over \(x\).  This proves both
adjunctions.  On a nonempty fiber, applying either fiberwise operation to a
pullback recovers the original value; this proves the two identity claims
when the map is surjective.  On a singleton fiber, pulling either fiberwise
operation back recovers \(w\); this proves the identity claims when the map
is injective.  Pointwise on the fiber through \(y\), the two source-side
composites take respectively the infimum and the supremum of the values of
\(w\) on that fiber.  They are therefore the claimed minorant and majorant;
equality for either composite is equivalent to constancy on every fiber,
which is already the pullback-image criterion.
\end{proof}

\begin{theorem}[Generic-norm trace descent for finite curve maps]
\label{thm:p126-curve-generic-norm-trace-descent}
Let
\(\varphi:(R,\mathfrak m)\to(S,\mathfrak n)\) be a finite local
homomorphism between one-dimensional reduced Nagata noetherian local rings,
let \(\mathfrak a\subset R\) be \(\mathfrak m\)-primary, and assume that the
branch map
\[
 \rho_{\mathfrak a}:V_{\mathfrak aS}\longrightarrow V_{\mathfrak a}
\]
is surjective.  Put \(D=R^\nu\).  For a base branch \(\mathfrak q\), write
\[
 K_{\mathfrak q}:=\operatorname{Frac}(D_{\mathfrak q}),\qquad
 T_{\mathfrak q}:=S\otimes_RK_{\mathfrak q},\qquad
 \delta_{\mathfrak q}:=\dim_{K_{\mathfrak q}}T_{\mathfrak q}.
\]
For every source branch \(\mathfrak q'\) over \(\mathfrak q\), let
\(e_{\mathfrak q'/\mathfrak q}\) and
\(f_{\mathfrak q'/\mathfrak q}\) denote its ramification index and residue
degree.  Then:
\begin{enumerate}
\item The algebra \(T_{\mathfrak q}\) is a nonzero product of finite field
      extensions of \(K_{\mathfrak q}\), and
      \[
       \delta_{\mathfrak q}
          =\sum_{\mathfrak q'\mid\mathfrak q}
             e_{\mathfrak q'/\mathfrak q}
             f_{\mathfrak q'/\mathfrak q}.
       \tag{\(\triangleleft\triangleleft\)}
       \label{eq:p126-curve-generic-degree}
      \]
\item On the continuous or completed ideal-psh classes define
      \[
       (\operatorname{Tr}_{\varphi,\mathfrak a}u)
          (v_{\mathfrak q}^{\mathfrak a})
       :=\sum_{\mathfrak q'\mid\mathfrak q}
          e_{\mathfrak q'/\mathfrak q}
          f_{\mathfrak q'/\mathfrak q}\,
          u(v_{\mathfrak q'}^{\mathfrak aS}),
      \]
      and put
      \[
       (\operatorname{Av}_{\varphi,\mathfrak a}^{\mathrm{gen}}u)
          (v_{\mathfrak q}^{\mathfrak a})
       :=\delta_{\mathfrak q}^{-1}
          (\operatorname{Tr}_{\varphi,\mathfrak a}u)
          (v_{\mathfrak q}^{\mathfrak a}).
       \tag{\(\triangleright\triangleright\)}
       \label{eq:p126-curve-generic-average}
      \]
      A positive sum containing \(-\infty\) is interpreted as
      \(-\infty\).  Both maps preserve order and decreasing limits, and
      \[
       \rho_{\mathfrak a,*}u
          \leq\operatorname{Av}_{\varphi,\mathfrak a}^{\mathrm{gen}}u
          \leq\rho_{\mathfrak a,!}u,
       \qquad
       \operatorname{Av}_{\varphi,\mathfrak a}^{\mathrm{gen}}
          \rho_{\mathfrak a}^*=\operatorname{id}.
       \tag{\(\blacktriangleright\blacktriangleright\)}
       \label{eq:p126-curve-generic-average-retraction}
      \]
      Conversely, if the branch map is not surjective, pullback admits no
      left inverse even as a map of sets.  Thus branch surjectivity is the
      exact condition for a finite curve-link pullback to admit a
      pushforward retraction.  Extending the normalized average by the same
      formula to finite real-valued obstacles \(\psi\), one has
      \[
       \operatorname{Av}^{\mathrm{gen}}_{\varphi,\mathfrak a}
          (P^{\mathrm{id}}_{\mathfrak aS}\psi)
       \leq
       P^{\mathrm{id}}_{\mathfrak a}
          (\operatorname{Av}^{\mathrm{gen}}_{\varphi,\mathfrak a}\psi).
       \tag{\(\blacktriangle\blacktriangle\)}
       \label{eq:p126-curve-generic-average-envelope}
      \]
      Equality holds exactly when no fiber contains two points \(y,z\) with
      \(\psi(y)<0<\psi(z)\).
\item On finite real-valued functions the normalized average is a positive
      linear contraction for the sup norm.  Hence
      \[
       \Pi_{\varphi,\mathfrak a}
          :=\rho_{\mathfrak a}^*
             \operatorname{Av}^{\mathrm{gen}}_{\varphi,\mathfrak a}
       \tag{\(\circledast\)}
       \label{eq:p126-curve-generic-barycentric-projection}
      \]
      is a positive norm-one idempotent projection onto the fiber-constant
      functions, which are exactly the pullback image.  Write
      \(\mathsf{Meas}_{\mathbf R}(V)\) for the finite signed measures on a
      finite set \(V\).  Dually, for
      \(\mu\in\mathsf{Meas}_{\mathbf R}(V_{\mathfrak a})\), define
      \[
       (\operatorname{Lift}_{\varphi,\mathfrak a}\mu)
          (\{v_{\mathfrak q'}^{\mathfrak aS}\})
       :=\frac{e_{\mathfrak q'/\mathfrak q}
                    f_{\mathfrak q'/\mathfrak q}}
               {\delta_{\mathfrak q}}\,
          \mu(\{v_{\mathfrak q}^{\mathfrak a}\}).
       \tag{\(\circledcirc\)}
       \label{eq:p126-curve-generic-measure-lift}
      \]
      The lift is positive and linear, is an isometry for the total-variation
      norm, and satisfies
      \((\rho_{\mathfrak a})_\#
          \operatorname{Lift}_{\varphi,\mathfrak a}\mu=\mu\).  It obeys
      the exact pairing identity
      \[
       \sum_{v\in V_{\mathfrak a}}
          (\operatorname{Av}^{\mathrm{gen}}_{\varphi,\mathfrak a}u)(v)
          \mu(\{v\})
       =\sum_{w\in V_{\mathfrak aS}}
          u(w)(\operatorname{Lift}_{\varphi,\mathfrak a}\mu)(\{w\}).
       \tag{\(\bullet\)}
       \label{eq:p126-curve-generic-function-measure-pairing}
      \]
      for every finite real-valued \(u\).  In particular, it sends
      probability measures to probability measures.  Moreover,
      \[
       Q_{\varphi,\mathfrak a}
         :=\operatorname{Lift}_{\varphi,\mathfrak a}
              (\rho_{\mathfrak a})_\#
       \tag{\(\circledcirc\circledast\)}
       \label{eq:p126-curve-generic-measure-projection}
      \]
      is a positive norm-one idempotent projection on
      \(\mathsf{Meas}_{\mathbf R}(V_{\mathfrak aS})\).  Its image is exactly
      the signed
      measures \(\nu\) whose restriction to every fiber has the canonical
      conditional weights:
      \[
       \nu(\{v_{\mathfrak q'}^{\mathfrak aS}\})
        =\frac{e_{\mathfrak q'/\mathfrak q}
                    f_{\mathfrak q'/\mathfrak q}}
                {\delta_{\mathfrak q}}
          ((\rho_{\mathfrak a})_\#\nu)
             (\{v_{\mathfrak q}^{\mathfrak a}\}).
       \tag{\(\bullet\circledast\)}
       \label{eq:p126-curve-generic-measure-projection-image}
      \]
      Under the pairing above, \(Q_{\varphi,\mathfrak a}\) is the Banach
      adjoint of \(\Pi_{\varphi,\mathfrak a}\).  More precisely, let
      \(\mu\) be a full-support probability measure on
      \(V_{\mathfrak a}\), and set
      \(\nu=\operatorname{Lift}_{\varphi,\mathfrak a}\mu\).  For every
      real-valued base function \(h\) and source function \(u\),
      \[
       \operatorname{Av}^{\mathrm{gen}}_{\varphi,\mathfrak a}
          ((\rho_{\mathfrak a}^*h)u)
        =h\operatorname{Av}^{\mathrm{gen}}_{\varphi,\mathfrak a}u,
       \qquad
       \Pi_{\varphi,\mathfrak a}((\rho_{\mathfrak a}^*h)u)
        =(\rho_{\mathfrak a}^*h)\Pi_{\varphi,\mathfrak a}u.
       \tag{\(\mathrm{CE}\)}
       \label{eq:p126-curve-generic-module-identity}
      \]
      Thus \(\Pi_{\varphi,\mathfrak a}\) is the conditional expectation
      onto the fiber-constant functions for the finite probability space
      \((V_{\mathfrak aS},\nu)\).  For every \(1\leq p\leq\infty\),
      \[
       \|\Pi_{\varphi,\mathfrak a}u\|_{L^p(\nu)}
          \leq \|u\|_{L^p(\nu)}.
       \tag{\(L^p\)}
       \label{eq:p126-curve-generic-lp-contraction}
      \]
      If \(1<p<\infty\), equality holds exactly when \(u\) is constant
      on every fiber.  For \(p=2\), the projection is the orthogonal
      projection onto \(\rho_{\mathfrak a}^*L^2(\mu)\), and one has the
      exact variance decomposition
      \[
       \begin{split}
       \|u\|_{L^2(\nu)}^2
        &=\|\Pi_{\varphi,\mathfrak a}u\|_{L^2(\nu)}^2
          +\|u-\Pi_{\varphi,\mathfrak a}u\|_{L^2(\nu)}^2,\\
       \|u-\Pi_{\varphi,\mathfrak a}u\|_{L^2(\nu)}^2
        &=\sum_{\mathfrak q}\mu(\{v_{\mathfrak q}^{\mathfrak a}\})
          \sum_{\mathfrak q'\mid\mathfrak q}
          \frac{e_{\mathfrak q'/\mathfrak q}
                    f_{\mathfrak q'/\mathfrak q}}
               {\delta_{\mathfrak q}}
          \left|u(v_{\mathfrak q'}^{\mathfrak aS})
             -(\operatorname{Av}^{\mathrm{gen}}_{\varphi,\mathfrak a}u)
                (v_{\mathfrak q}^{\mathfrak a})\right|^2.
       \end{split}
       \tag{\(\mathrm{Var}\)}
       \label{eq:p126-curve-generic-variance}
      \]
\item The trace and average commute with the change-of-normalizing-ideal
      transports of
      \eqref{eq:p122-curve-psh-normalization-transport}.  In a tower of
      branch-surjective finite local maps, the unnormalized traces compose
      exactly.  As operators on all finite real-valued functions, the
      normalized averages compose if and only if the upper generic degree is
      constant on the intermediate branches above each fixed base branch.
      Dually, the signed-measure lifts compose under exactly the same
      condition.  In particular, both compose when the upper map is finite
      flat.
\item Let \(g\in\mathfrak n\) be an \(S\)-nonzerodivisor, and set
      \[
       \beta_{\mathfrak q}(g)
          :=\sum_{\mathfrak q'\mid\mathfrak q}
             f_{\mathfrak q'/\mathfrak q}\,
             w_{\mathfrak q'}(g).
       \tag{\(\blacktriangleleft\blacktriangleleft\)}
       \label{eq:p126-curve-generic-norm-order}
      \]
      This is the order of the generic determinant norm of \(g\) over the
      normalized base DVR.  For every sufficiently large integer \(M\), and
      every sufficiently large common multiple \(L\) of the positive
      integers \(\delta_{\mathfrak q}\), there are
      \(R\)-nonzerodivisors \(h_M,\widetilde h_L\in\mathfrak m\) with
      \[
       w_{\mathfrak q}(h_M)=M\beta_{\mathfrak q}(g),
       \qquad
       w_{\mathfrak q}(\widetilde h_L)
          =\frac{L}{\delta_{\mathfrak q}}\beta_{\mathfrak q}(g),
      \]
      and, for every \(N\geq1\),
      \[
       \operatorname{Tr}_{\varphi,\mathfrak a}
          u^{\mathfrak aS}_{g,N}
          =u^{\mathfrak a}_{h_M,MN},
       \qquad
       \operatorname{Av}^{\mathrm{gen}}_{\varphi,\mathfrak a}
          u^{\mathfrak aS}_{g,N}
          =u^{\mathfrak a}_{\widetilde h_L,LN}.
       \tag{\(\blacklozenge\blacklozenge\)}
       \label{eq:p126-curve-generic-principal-descent}
      \]
\end{enumerate}
Thus every branch-surjective finite curve map has a canonical generic-degree
weighted link pushforward, and its principal source models descend exactly
after a conductor-controlled common scaling.  The next theorem identifies
the finite-flat range in which the generic norm is already one canonical
downstairs determinant norm, with no conductor scaling.  The signed-measure
lift and conditional expectation above are only finite-link operations;
they are not a construction or functoriality assertion for a singular
Monge--Amp\`ere measure.
\end{theorem}

\begin{proof}
Fix \(\mathfrak q\), let \(\mathfrak p\) be its corresponding minimal prime
of \(R\), and put \(A=D_{\mathfrak q}\).  Since
\(K_{\mathfrak q}\simeq R_{\mathfrak p}\), the algebra
\(T_{\mathfrak q}=S_{R\setminus\mathfrak p}\) is reduced and finite over
\(K_{\mathfrak q}\).  Branch surjectivity makes it nonzero, so it is a
product of finite field extensions.  The DVR \(A\) is Nagata.  Hence its
integral closure \(B\) in \(T_{\mathfrak q}\) is finite and torsion-free,
and therefore free of rank \(\delta_{\mathfrak q}\).  Its maximal ideals
are precisely the source branches above \(\mathfrak q\).  If \(\pi\) is a
uniformizer of \(A\), filtering the local factors of \(B/\pi B\) gives
\[
 \delta_{\mathfrak q}=\operatorname{length}_A(B/\pi B)
   =\sum_{\mathfrak q'\mid\mathfrak q}
      e_{\mathfrak q'/\mathfrak q}
      f_{\mathfrak q'/\mathfrak q},
\]
which proves \eqref{eq:p126-curve-generic-degree}.

The coefficients of the normalized average are positive and sum to one.
The two inequalities in
\eqref{eq:p126-curve-generic-average-retraction} are therefore the elementary
minimum--average--maximum inequalities on each fiber.  By
\eqref{eq:p126-ramification-cancellation}, a pullback is constant on every
fiber, so the degree formula gives the retraction identity.  The remaining
order and limit assertions are coordinatewise on a finite set.  If the
branch map is not surjective, \(\rho_{\mathfrak a}^*\) is not injective by
\Cref{thm:p122-curve-ideal-psh-functoriality}; a map with a left inverse is
injective, proving necessity.  The function \(s\mapsto\min\{s,0\}\) is
concave.  Its finite weighted Jensen inequality on each branch fiber gives
\eqref{eq:p126-curve-generic-average-envelope}.  Since every weight is
strictly positive, equality in this piecewise-linear Jensen inequality fails
precisely when the fiber contains both a strictly negative and a strictly
positive value, which proves the stated criterion.

For finite real-valued functions, a positive weighted average has sup norm
at most the source sup norm.  Together with
\(\operatorname{Av}^{\mathrm{gen}}\rho^*=\operatorname{id}\), this makes
\(\Pi_{\varphi,\mathfrak a}\) idempotent with the fiber-constant image; its
norm is one because it fixes constants.  Formula
\eqref{eq:p126-curve-generic-degree} says that the masses assigned in
\eqref{eq:p126-curve-generic-measure-lift} sum over each fiber to the original
base mass.  This proves both the probability and pushforward identities.
Expanding the two finite sums proves
\eqref{eq:p126-curve-generic-function-measure-pairing}.
All lifted atoms over one base point have the sign of that base mass, so
summing absolute values fiber by fiber gives
\[
 \|\operatorname{Lift}_{\varphi,\mathfrak a}\mu\|_{\mathrm{TV}}
   =\|\mu\|_{\mathrm{TV}}.
\]
Together with the evident positivity and linearity, this proves the stated
total-variation isometry on signed measures.  Since
\((\rho_{\mathfrak a})_\#\operatorname{Lift}_{\varphi,\mathfrak a}
=\operatorname{id}\), the operator \(Q_{\varphi,\mathfrak a}\) is
idempotent and has precisely the image described in
\eqref{eq:p126-curve-generic-measure-projection-image}.  Measure pushforward
is a total-variation contraction, while the lift is an isometry; hence
\(\|Q_{\varphi,\mathfrak a}\|\leq1\), and equality follows because \(Q\)
fixes its nonzero image.  The pairing identity applied to
\((\rho_{\mathfrak a})_\#\nu\) says exactly that \(Q\) is the adjoint of
\(\Pi\).

The first identity in
\eqref{eq:p126-curve-generic-module-identity} follows by taking the base
factor \(h\) outside each weighted fiber sum, and the second follows after
pullback.  With \(\nu=\operatorname{Lift}\mu\), Jensen's inequality for
\(t\mapsto |t|^p\), applied on every fiber and then integrated against
\(\mu\), proves \eqref{eq:p126-curve-generic-lp-contraction} for finite
\(p\); the case \(p=\infty\) was proved above.  Since all fiber weights and
all base masses are positive, strict convexity gives the stated equality
criterion when \(1<p<\infty\).  If \(z=\rho_{\mathfrak a}^*h\) is
fiber-constant, the pairing identity and the retraction formula give
\[
 \langle u-\Pi_{\varphi,\mathfrak a}u,z\rangle_{L^2(\nu)}
 =\left\langle
    \operatorname{Av}^{\mathrm{gen}}_{\varphi,\mathfrak a}
       (u-\Pi_{\varphi,\mathfrak a}u),h
  \right\rangle_{L^2(\mu)}=0.
\]
Thus \(\Pi_{\varphi,\mathfrak a}\) is the asserted orthogonal projection.
Pythagoras gives the first line of
\eqref{eq:p126-curve-generic-variance}, and expanding the lifted measure
fiber by fiber gives its second line.

Changing the normalizing ideal multiplies all source coordinates over one
base branch by the same positive factor, which proves transport
compatibility.  In a composable tower, ramification indices multiply by
\cite[Tag~0BRL]{StacksProject}, and residue degrees multiply by the field
tower law.  Expanding the two finite sums proves exact composition of the
unnormalized traces.  If the upper degree has the common value \(c_{\mathfrak
q}\) over the intermediate branches above \(\mathfrak q\), the composite
degree is \(c_{\mathfrak q}\delta_{\mathfrak q}\); division proves the
stated composition law for normalized averages.
Conversely, let \(D_{\mathfrak q}\) be the composite generic degree and let
\(c_{\mathfrak q'}\) be the upper generic degree at an intermediate branch.
Compare the coefficient of a top branch over \(\mathfrak q'\) in the direct
and iterated normalized averages.  After cancelling the positive product of
ramification and residue degrees, equality of the two operators forces
\[
                   \delta_{\mathfrak q}c_{\mathfrak q'}=D_{\mathfrak q}
\]
for every intermediate branch over \(\mathfrak q\).  The right side is
fixed, so the upper generic degree is constant over that intermediate fiber.
The lift statement follows either from the same coefficient comparison or from
the exact function--measure duality.

Finally, \(g\) is nonzero in every field factor of
\(T_{\mathfrak q}\).  Multiplication by \(g\) on \(B\) is injective, and
the DVR length--determinant formula
\cite[\emph{Lengths and determinants}, Tag~02QA]{stacks-project} gives
\[
 \operatorname{ord}_A\det(m_g\mid B)
   =\operatorname{length}_A(B/gB)
   =\sum_{\mathfrak q'\mid\mathfrak q}
      f_{\mathfrak q'/\mathfrak q}w_{\mathfrak q'}(g)
   =\beta_{\mathfrak q}(g).
\]
Every \(\beta_{\mathfrak q}(g)\) is a positive integer.  Let
\(C_{\mathfrak q}\) be the conductor thresholds of
\eqref{eq:p126-curve-conductor-orders}.  For all sufficiently large \(M\),
the integers \(M\beta_{\mathfrak q}(g)\) exceed those thresholds.  For all
sufficiently large common multiples \(L\) of the
\(\delta_{\mathfrak q}\), the same holds for
\(L\beta_{\mathfrak q}(g)/\delta_{\mathfrak q}\).  Applying
\Cref{prop:p126-curve-conductor-principal-realization} to these two vectors
produces \(h_M\) and \(\widetilde h_L\).  Since
\(d_{\mathfrak q'}(\mathfrak aS)
 =e_{\mathfrak q'/\mathfrak q}d_{\mathfrak q}(\mathfrak a)\), direct
substitution gives
\[
 (\operatorname{Tr}_{\varphi,\mathfrak a}
       u^{\mathfrak aS}_{g,N})(v_{\mathfrak q}^{\mathfrak a})
   =-\frac{\beta_{\mathfrak q}(g)}
           {N d_{\mathfrak q}(\mathfrak a)},
\]
and division by \(\delta_{\mathfrak q}\) gives the average.  The prescribed
orders of \(h_M\) and \(\widetilde h_L\) now yield exactly
\eqref{eq:p126-curve-generic-principal-descent}.
\end{proof}

\begin{example}[Variable generic degree obstructs normalized composition]
\label{ex:p126-curve-variable-degree-average}
Let \(K\) be a field and consider the finite local maps
\[
 \begin{aligned}
 R&=K[[t]]\longrightarrow
 S=K[[x,y]]/(xy),& t&\longmapsto x+y,\\
 S&\longrightarrow
 U=K[[a,b,c]]/(ab,ac,bc),& x&\longmapsto a,\quad
 y\longmapsto b+c.
 \end{aligned}
\]
Indeed, \(x\) satisfies \(x^2-tx=0\), while \(b\) satisfies
\(b^2-yb=0\) and \(c=y-b\).  Thus the first map has two generic branches
of degree one.  Over its \(x\)-branch the second map has one generic branch,
whereas over its \(y\)-branch it has two, again all with ramification and
residue degree one.

Take \(\mathfrak a=(t)\), and give the \(a,b,c\) branches of the completed
link of \(U\) the values \(-1,-2,-3\), respectively.  The direct normalized
average from \(U\) to \(R\) is
\[
                          \frac{-1-2-3}{3}=-2.
\]
Averaging first to \(S\) gives the pair \((-1,-5/2)\), whose subsequent
average to \(R\) is
\[
                          \frac{-1-5/2}{2}=-\frac74.
\]
Thus normalized generic averages need not compose when the upper generic
degree varies between intermediate branches.  The unnormalized traces do
compose: both routes give \(-6\), exactly as asserted in
\Cref{thm:p126-curve-generic-norm-trace-descent}.  The dual obstruction is
equally explicit.  A unit mass on the unique branch of \(R\) lifts directly
to masses
\[
                         (1/3,1/3,1/3)
\]
on the \(a,b,c\) branches, whereas successive lifts through \(S\) give
\[
                         (1/2,1/4,1/4).
\]
Thus the signed-measure lifts also fail to compose, with the same sharp
constant-upper-degree boundary.
\end{example}

\begin{theorem}[Finite-flat norm descent on reduced curve links]
\label{thm:p126-curve-finite-flat-norm-descent}
Let
\(\varphi:(R,\mathfrak m)\to(S,\mathfrak n)\) be a finite flat local
homomorphism of constant rank \(d\geq1\) between one-dimensional reduced
Nagata noetherian local rings, and let \(\mathfrak a\subset R\) be
\(\mathfrak m\)-primary.  Put \(D=R^\nu\) and \(E=S^\nu\), and write
\(v_{\mathfrak q}^{\mathfrak a}
=w_{\mathfrak q}/d_{\mathfrak q}(\mathfrak a)\) and
\(v_{\mathfrak q'}^{\mathfrak aS}
=w_{\mathfrak q'}/d_{\mathfrak q'}(\mathfrak aS)\).
For a branch \(\mathfrak q'\) of \(S\) over a branch
\(\mathfrak q\) of \(R\), let
\[
 e_{\mathfrak q'/\mathfrak q}\geq1,
 \qquad
 f_{\mathfrak q'/\mathfrak q}
   :=[\kappa(\mathfrak q'): \kappa(\mathfrak q)]
\]
be its ramification index and residue degree.  Then:
\begin{enumerate}
\item The branch map \(\rho_{\mathfrak a}\) is surjective, the generic
      degrees of
      \Cref{thm:p126-curve-generic-norm-trace-descent} all equal \(d\), and
      over every \(\mathfrak q\) one has
      \[
       \sum_{\mathfrak q'\mid\mathfrak q}
          e_{\mathfrak q'/\mathfrak q}
          f_{\mathfrak q'/\mathfrak q}=d.
       \tag{\(\triangleleft\)}\label{eq:p126-curve-ef-degree}
      \]
\item For every nonzerodivisor \(g\in S\), its finite-flat determinant
      norm satisfies
      \[
       w_{\mathfrak q}\!\left(\mathrm N_{S/R}(g)\right)
          =\sum_{\mathfrak q'\mid\mathfrak q}
             f_{\mathfrak q'/\mathfrak q}\,
             w_{\mathfrak q'}(g).
       \tag{\(\blacktriangle\)}\label{eq:p126-curve-norm-order}
      \]
      If \(g\in\mathfrak n\), then \(\mathrm N_{S/R}(g)\) is an
      \(R\)-nonzerodivisor in \(\mathfrak m\), so its principal ideal is
      \(\mathfrak m\)-primary.
\item On the continuous or completed ideal-psh classes define the
      finite-flat norm average by
      \[
       (\operatorname{Av}_{\varphi,\mathfrak a}u)
          (v_{\mathfrak q}^{\mathfrak a})
       :=\frac1d\sum_{\mathfrak q'\mid\mathfrak q}
          e_{\mathfrak q'/\mathfrak q}
          f_{\mathfrak q'/\mathfrak q}\,
          u(v_{\mathfrak q'}^{\mathfrak aS}),
      \tag{\(\triangleright\)}\label{eq:p126-curve-norm-average}
      \]
      This is exactly the generic-degree average
      \eqref{eq:p126-curve-generic-average} in the finite-flat case; a
      positive-weighted sum containing \(-\infty\) is interpreted as
      \(-\infty\).  This is order preserving, preserves decreasing limits,
      and satisfies
      \[
       \rho_{\mathfrak a,*}u
          \leq\operatorname{Av}_{\varphi,\mathfrak a}u
          \leq\rho_{\mathfrak a,!}u,
       \qquad
       \operatorname{Av}_{\varphi,\mathfrak a}
          \rho_{\mathfrak a}^*=\operatorname{id}.
       \tag{\(\blacktriangleright\)}\label{eq:p126-curve-average-sandwich}
      \]
      It commutes with the change-of-normalizing-ideal transports
      \eqref{eq:p122-curve-psh-normalization-transport}, and the norm
      averages compose in towers of finite flat local maps.
\item Norm averaging sends principal models to principal models exactly:
      for every \(S\)-nonzerodivisor \(g\in\mathfrak n\) and \(N\geq1\),
      \[
       \operatorname{Av}_{\varphi,\mathfrak a}
          u^{\mathfrak aS}_{g,N}
        =u^{\mathfrak a}_{\mathrm N_{S/R}(g),dN}.
       \tag{\(\blacklozenge\)}\label{eq:p126-curve-principal-norm-descent}
      \]
\end{enumerate}
Consequently, the source conductor realization of every sufficiently deep
integral branch-order vector descends by one determinant norm to the
weighted-average principal profile downstairs.  This supplies a genuine
pushforward and conductor-compatible range of P12.6 for finite flat curve
maps.  It asserts no norm pushforward for a nonflat finite map.
\end{theorem}

\begin{proof}
Because \(R\) is local and \(S\) is finite flat of constant rank \(d\),
the \(R\)-module \(S\) is free of rank \(d\), and
\(\mathrm N_{S/R}(g)=\det(m_g)\) is the determinant norm of
\eqref{eq:appA-element-norm}.

Fix \(\mathfrak q\), let \(A=D_{\mathfrak q}\) be the corresponding
normalized discrete valuation ring of \(R\), and write \(K\) for its
fraction field.  If \(\mathfrak p\) is the minimal prime of \(R\)
corresponding to \(\mathfrak q\), then \(K\simeq R_{\mathfrak p}\).
Consequently the generic algebra
\[
                          T:=S\otimes_RK
\]
is the localization \(S_{R\setminus\mathfrak p}\).  It is therefore a
reduced \(K\)-algebra of dimension \(d\), hence a product of finite field
extensions of \(K\).  The ring \(A\) is Nagata, being a localization of the
finite normalization of the Nagata ring \(R\).  Thus the integral closure
\(B\) of \(A\) in \(T\) is finite; it is torsion-free, hence a finite free
\(A\)-module of rank \(d\).  Its maximal ideals are exactly the branches
\(\mathfrak q'\) above \(\mathfrak q\), and
\(B_{\mathfrak q'}\) is their normalized discrete valuation ring.

Let \(\pi\) be a uniformizer of \(A\).  The quotient \(B/\pi B\) has
\(A\)-length \(d\).  Localizing it at the maximal ideals of \(B\), and
filtering each
\(B_{\mathfrak q'}/\pi B_{\mathfrak q'}
 =B_{\mathfrak q'}/\mathfrak q'^{e_{\mathfrak q'/\mathfrak q}}\),
gives
\[
 d=\operatorname{length}_A(B/\pi B)
   =\sum_{\mathfrak q'\mid\mathfrak q}
      e_{\mathfrak q'/\mathfrak q}
      [\kappa(\mathfrak q'): \kappa(\mathfrak q)].
\]
This proves \eqref{eq:p126-curve-ef-degree}, and in particular proves that
every base branch has a branch above it.

If \(g\) is a nonzerodivisor of \(S\), it is nonzero in every field factor
of \(T\), so multiplication by \(g\) is injective on \(B\).  Determinants
commute with base change by \Cref{prop:appA-norm-identities}.  After
extending scalars to \(K\), multiplication on \(S\) and on the
\(A\)-lattice \(B\) becomes the same endomorphism of \(T\); its determinant
is independent of the chosen \(K\)-basis.  Hence the determinant of
\(m_g\) on \(B\) is the image of \(\mathrm N_{S/R}(g)\) in \(A\).  The DVR length--determinant formula
\cite[\emph{Lengths and determinants}, Tag~02QA]{stacks-project} and the
same residue-field filtration give
\[
 \begin{aligned}
 w_{\mathfrak q}\!\left(\mathrm N_{S/R}(g)\right)
   &=\operatorname{length}_A(B/gB)\\
   &=\sum_{\mathfrak q'\mid\mathfrak q}
       f_{\mathfrak q'/\mathfrak q}\,
       \operatorname{length}_{B_{\mathfrak q'}}
          (B_{\mathfrak q'}/gB_{\mathfrak q'})\\
   &=\sum_{\mathfrak q'\mid\mathfrak q}
       f_{\mathfrak q'/\mathfrak q}\,
       w_{\mathfrak q'}(g),
 \end{aligned}
\]
which is \eqref{eq:p126-curve-norm-order}.  It is finite on every branch,
so the norm avoids every minimal prime of the reduced ring \(R\) and is a
nonzerodivisor.  If \(g\in\mathfrak n\), its image in the artinian local
algebra \(S/\mathfrak mS\) is nilpotent.  Multiplication by that image has
zero determinant, so \(\mathrm N_{S/R}(g)\in\mathfrak m\).  The resulting
principal ideal is \(\mathfrak m\)-primary by the one-dimensional argument
in \Cref{prop:p126-curve-conductor-principal-realization}.

Formula \eqref{eq:p126-curve-ef-degree} says that the coefficients in
\eqref{eq:p126-curve-norm-average} are positive and sum to one on each
fiber.  Thus the average lies between the fiberwise infimum and supremum,
which proves the first part of
\eqref{eq:p126-curve-average-sandwich}; it plainly preserves finite values,
order, and decreasing limits.  The ramification-cancellation formula
\eqref{eq:p126-ramification-cancellation} makes a pullback constant on every
fiber, so averaging and using \eqref{eq:p126-curve-ef-degree} proves the
retraction identity.  A change of normalizing ideal multiplies all source
coordinates over a fixed base branch by the same positive factor, proving
compatibility with \eqref{eq:p122-curve-psh-normalization-transport}.
Ramification indices and residue degrees multiply in a tower; expanding the
two successive finite sums therefore proves composition of the averages.

Finally,
\[
 \begin{aligned}
 (\operatorname{Av}_{\varphi,\mathfrak a}
       u^{\mathfrak aS}_{g,N})(v_{\mathfrak q}^{\mathfrak a})
  &=-\frac1d\sum_{\mathfrak q'\mid\mathfrak q}
       e_{\mathfrak q'/\mathfrak q}f_{\mathfrak q'/\mathfrak q}
       \frac{w_{\mathfrak q'}(g)}
            {N e_{\mathfrak q'/\mathfrak q}
               d_{\mathfrak q}(\mathfrak a)}\\
  &=-\frac{w_{\mathfrak q}(\mathrm N_{S/R}(g))}
            {dN d_{\mathfrak q}(\mathfrak a)}
   =u^{\mathfrak a}_{\mathrm N_{S/R}(g),dN}
       (v_{\mathfrak q}^{\mathfrak a}),
 \end{aligned}
\]
where the middle equality is \eqref{eq:p126-curve-norm-order}.  This proves
\eqref{eq:p126-curve-principal-norm-descent}; combining it with the source
version of \Cref{prop:p126-curve-conductor-principal-realization} gives the
last assertion.
\end{proof}

\begin{example}[The two endpoint weights in finite-flat norm descent]
\label{ex:p126-curve-norm-endpoint-weights}
Let \(K\) be a field.  For the totally ramified map
\[
 K[[t]]\longrightarrow K[[u]],\qquad t\longmapsto u^e,
\]
the rank is \(d=e\), the residue degree is one, and
\(\mathrm N(u)=\pm t\).  With \(\mathfrak a=(t)\), the normalized source
value of \(u^{(u^e)}_{u,N}\) is \(-1/(eN)\); the weight \(ef/d=1\) leaves
it unchanged, and the descended model
\(u^{(t)}_{\pm t,eN}\) has the same value.  Thus the ramification index
belongs in the averaging weight but not in
\eqref{eq:p126-curve-norm-order}.

At the other endpoint, if \(K'/K\) is a finite extension of degree \(f\),
then
\[
 K[[t]]\longrightarrow K'[[t]]
\]
has \(e=1\), \(d=f\), and \(\mathrm N(t)=t^f\).  The source model
\(u^{(t)}_{t,N}\) and the descended model
\(u^{(t)}_{t^f,fN}\) both have value \(-1/N\).  This time the residue
degree appears both in the norm order and in the rank denominator and
cancels there.
\end{example}

\begin{example}[A boundary profile on the two-branch node]
\label{ex:p122-curve-ideal-psh-boundary}
For \(R=\mathbf C[[x,y]]/(xy)\) and
\(\mathfrak a=(x,y)\), label the two points of
\(V_{\mathfrak a}\) by the \(x\)- and \(y\)-axes.  The profile
\[
                             u=(-1,0)
\]
belongs to \(\operatorname{CPSH}^{\mathrm{id}}_{\mathfrak a}(R)\) and
is an increasing uniform limit of principal-ideal profiles.  It is not a
decreasing limit of principal-ideal profiles: every such profile is
strictly negative on both branches.  Thus uniform ideal closure and
decreasing model regularization genuinely differ at a reducible boundary,
even though the completed psh envelope is the coordinatewise formula
\eqref{eq:p123-curve-ideal-psh-envelope}.
\end{example}

\begin{example}[A nonsurjective branch pullback reaches the bottom]
\label{ex:p122-curve-psh-pullback-bottom}
Let
\[
 R=\mathbf C[[x,y]]/(xy),\qquad S=\mathbf C[[x]],
 \qquad
 \varphi:R\longrightarrow S,\qquad x\longmapsto x,\quad y\longmapsto0,
\]
and take \(\mathfrak a=(x,y)\), so that
\(\mathfrak aS=(x)\).  This is a finite local homomorphism of reduced
excellent curve rings.  Its normalized-link map sends the unique branch of
\(S\) to the \(x\)-axis branch of \(R\), and misses the \(y\)-axis branch.
On the two-point link of \(R\), ordered as the \(x\)- and \(y\)-axes, the
function
\[
                              u=(-\infty,0)
\]
belongs to \(\operatorname{PSH}^{\mathrm{id},\bot}_{\mathfrak a}(R)\)
but is not its bottom element.  Nevertheless,
\[
                 \rho_{\mathfrak a}^*u\equiv-\infty
                 \quad\text{on }V_{(x)}(S).
\]
Thus finite pullback preserves the bottom-inclusive completion but cannot in
general preserve the class obtained by deleting its bottom element.  This is
the precise reason for retaining \(-\infty\) in the functorial statement.
Here \(\rho_{\mathfrak a}\) is injective but not surjective, so pullback is
surjective but not injective, exactly as the preceding theorem predicts.
\end{example}

The word ``psh'' in the preceding notation cannot be justified by extending
every homogeneous branch profile radially across an analytic singularity.
On any isolated curve singularity whose normalization is a finite union of
discs and whose branches are separated by the conductor, however, the exact
branch-compatibility condition can be computed.

\begin{theorem}[Conductorial radial classification on a disc-normalized curve]
\label{thm:p122-conductorial-disc-radial-psh-classification}
Let \(K\) be a complete nontrivially valued non-Archimedean field, choose
\(\varpi\in K\) with \(0<|\varpi|<1\), and put
\(r=|\varpi|\) and \(a=-\log r\).  Let \(A\) be a reduced strict
\(K\)-affinoid algebra of pure dimension one whose normalization is finite
and has the form
\[
 A\lhook\joinrel\longrightarrow
 B=\prod_{i=1}^q K\langle t_i\rangle,
 \qquad
 \nu:\widetilde U:=\mathcal M(B)\longrightarrow U:=\mathcal M(A).
\]
Suppose that there is a \(K\)-rational point \(o\in U\) such that
\(\nu^{-1}(o)=\{p_1,\ldots,p_q\}\), where \(p_i\) is the origin of the
\(i\)-th disc \(D_i:=\mathcal M(K\langle t_i\rangle)\), and that \(\nu\)
is an isomorphism away from these points.  Assume that the conductor
contains a separated branch ideal
\[
 \mathfrak d_{\boldsymbol e}
   :=\prod_{i=1}^q t_i^{e_i}K\langle t_i\rangle
   \ \subseteq\
 \mathfrak c_{B/A}:=\{b\in B:bB\subseteq A\}
 \tag{\(\blacktriangledown\)}
 \label{eq:p122-disc-conductor-branch-ideal}
\]
for some integers \(e_i\geq1\).  This conductor hypothesis is automatic if
the finite \(A\)-module \(B/A\) is supported at the maximal ideal of \(o\).

For \(c=(c_1,\ldots,c_q)\in[-\infty,0]^q\), define
\(h_c^\circ:U\setminus\{o\}\to[-\infty,0]\) through the normalization by
\[
 h_c^\circ(\nu(z)):=
 \begin{cases}
  -\infty,&c_i=-\infty,\\[1mm]
  -\dfrac{c_i}{a}\log|t_i(z)|,&c_i> -\infty,
 \end{cases}
 \qquad z\in D_i\setminus\{p_i\}.
 \tag{\(\blacktriangleleft\)}
 \label{eq:p122-disc-normalized-radial-function}
\]
Then the following are equivalent:
\begin{enumerate}
\item \(h_c^\circ\) extends to a classically psh function on \(U\);
\item \(h_c^\circ\) extends to a CLD-v2 psh function on \(U\);
\item either \(c=0\), or \(c_i<0\) for every \(i\), with
      \(-\infty<0\) allowed.
\end{enumerate}
The extension is unique.  Its value at \(o\) is zero for \(c=0\), and is
\(-\infty\) otherwise.  Thus the exact radial-extendable cone is
\[
                         \{0\}\cup[-\infty,0)^q.
 \tag{\(\blacktriangleright\)}
 \label{eq:p122-disc-normalized-radial-cone}
\]
Every finite strictly negative rational profile is realized by one scaled
logarithm \((aN)^{-1}\log|f|\) with \(f\in A\), and all-negative extended
profiles are obtained by decreasing limits after putting zero in the
normalization components corresponding to their \(-\infty\) coordinates.
\end{theorem}

\begin{proof}
First note why the final sentence of the conductor hypothesis holds.  Let
\(\mathfrak m\subset A\) be the maximal ideal of \(o\).  If
\(\operatorname{Supp}_A(B/A)\subseteq\{\mathfrak m\}\), then noetherianity
gives an integer \(M\geq1\) with
\(\mathfrak m^M(B/A)=0\), hence
\[
                            \mathfrak m^M B\subseteq A.
\]
Writing \(I_i\) for the image of \(\mathfrak mB\) in
\(K\langle t_i\rangle\), the fiber hypothesis and the affinoid
Nullstellensatz give \(\sqrt{I_i}=(t_i)\).  Choose \(d_i\geq1\) with
\(t_i^{d_i}\in I_i\), and set \(e_i=Md_i\).  Then
\[
 \prod_i t_i^{e_i}K\langle t_i\rangle
       \subseteq (\mathfrak mB)^M=\mathfrak m^M B\subseteq A.
\]
The left-hand side is a \(B\)-ideal, so it is contained in
\(\mathfrak c_{B/A}\), proving
\eqref{eq:p122-disc-conductor-branch-ideal}.

For the psh classification, the zero profile has the constant zero
extension.  Suppose next that every \(c_i<0\), and put
\[
 I:=\{i:c_i> -\infty\},\qquad J:=\{i:c_i=-\infty\}.
\]
If \(I=\varnothing\), use the constant \(-\infty\) extension.  Otherwise,
put \(\lambda_i=-c_i\in(0,+\infty)\) for \(i\in I\).  Choose positive
rationals
\[
                  \lambda_i^{(n)}=k_{n,i}/N_n\nearrow\lambda_i
                  \qquad(i\in I)
\]
with one denominator \(N_n\) at each stage and with
\(k_{n,i}\geq e_i\).  The latter condition costs nothing: after choosing
common-denominator rational approximants, multiply all their numerators and
their common denominator by one sufficiently large integer.  Define
\(b_n=(b_{n,1},\ldots,b_{n,q})\in B\) by
\[
 b_{n,i}:=
 \begin{cases}
   t_i^{k_{n,i}},&i\in I,\\
   0,&i\in J.
 \end{cases}
\]
By \eqref{eq:p122-disc-conductor-branch-ideal}, \(b_n\in A\); denote the
corresponding analytic function on \(U\) by \(f_n\).  Then
\(\nu^*f_n=b_n\), and therefore
\[
 g_n:=(aN_n)^{-1}\log|f_n|
\]
pulls back to
\(a^{-1}\lambda_i^{(n)}\log|t_i|\) on \(D_i\) for \(i\in I\), and to
\(-\infty\) on \(D_j\) for \(j\in J\).  It also takes the value
\(-\infty\) at \(o\).  Each \(g_n\) is CLD-v2 psh by
\Cref{prop:ch15-log-approachable,prop:ch15-psh-cone}.
\Cref{prop:ch15-log-approachable} applies to every holomorphic function, so
this remains valid when \(J\ne\varnothing\) and \(f_n\) is a zero-divisor;
the proposition's additional non-zero-divisor hypothesis concerns only local
integrability.  Since \(\log|t_i|\leq0\), the sequence decreases pointwise
to the claimed
extension of \(h_c^\circ\), which is CLD-v2 psh by decreasing closedness.
This proves (3)\(\Rightarrow\)(2), while (2)\(\Rightarrow\)(1) follows from
\Cref{thm:cld-to-classical-psh}.

Conversely, let \(h\) be a classically psh extension and suppose that
\(c_i=0\) for some \(i\).  Its pullback \(\widetilde h=h\circ\nu\) is
classically psh by \cite[Proposition~4.2(6)]{GR25Classical}.  The points
\(t_i=\varpi^m\) tend to \(p_i\), so upper
semicontinuity gives \(h(o)=\widetilde h(p_i)\geq0\).  If \(c\ne0\), choose
\(j\) with \(c_j<0\).  On the closed radius-\(r\) subdisc of \(D_j\), the
unique Shilov point has value \(c_j<0\), and every other noncentral value is
at most \(c_j\) (or is \(-\infty\) when \(c_j=-\infty\)).  The affinoid
maximum principle \cite[Theorem~6.4]{GR25Classical} contradicts
\(\widetilde h(p_j)=h(o)\geq0\).  This proves (1)\(\Rightarrow\)(3).

For uniqueness, the zero branch limit gives \(h(o)\geq0\) when \(c=0\),
and the maximum principle on any branch disc gives \(h(o)\leq0\).  In the
all-negative case, a finite coordinate \(c_j\) and the radius-\(r^m\) disc
give \(h(o)\leq m c_j\) for every \(m\), hence \(h(o)=-\infty\); if every
coordinate is \(-\infty\), one branch disc gives the same conclusion.
Finally, a finite rational vector \(c_i=-k_i/N\) can be represented with
all \(k_i\geq e_i\) by scaling numerator and denominator, and the preceding
conductor construction realizes it by one \((aN)^{-1}\log|f|\).  The same
coordinatewise rational approximation proves the asserted decreasing
closure on the radius-\(r\) branch link.
\end{proof}

\begin{example}[The one-branch cusp]
\label{ex:p122-cusp-radial-psh-classification}
Let
\[
 A_{\mathrm{cusp}}:=K\langle u,v\rangle/(v^2-u^3)
      \lhook\joinrel\longrightarrow B:=K\langle t\rangle,
 \qquad u\longmapsto t^2,\quad v\longmapsto t^3.
\]
The decomposition
\(B=K\langle t^2\rangle\oplus tK\langle t^2\rangle\) makes this extension
finite, the element \(t\) is integral because \(t^2=u\), and it equals
\(v/u\) in the fraction field.  Since \(B\) is normal, this is the finite
normalization.  Inside \(B\) one has
\[
 A_{\mathrm{cusp}}=K+t^2K\langle t\rangle,
 \qquad
 \mathfrak c_{B/A_{\mathrm{cusp}}}=t^2K\langle t\rangle.
\]
Indeed, even--odd decomposition gives
\[
 t^2K\langle t\rangle
   =t^2K\langle t^2\rangle\oplus t^3K\langle t^2\rangle
   \subseteq A_{\mathrm{cusp}},
\]
whereas every nonconstant monomial in the generators \(t^2,t^3\) has order
at least two; this proves the first equality.  The displayed ideal is
therefore contained in the conductor.  Conversely, if
\(b=\gamma+t^2g(t)\) lies in the conductor, then
\(tb\in A_{\mathrm{cusp}}\); the coefficient of \(t\) forces
\(\gamma=0\), proving the reverse inclusion.
The normalization is an isomorphism away from the origin, so
\Cref{thm:p122-conductorial-disc-radial-psh-classification} applies with
\(q=1\) and \(e_1=2\).  Hence every
\(c\in[-\infty,0]=\{0\}\cup[-\infty,0)\) has a unique radial extension
which is both CLD-v2 and classically psh.  Its value at the cusp is zero for
\(c=0\) and \(-\infty\) otherwise.  More explicitly, if
\(c=-p/s<0\) with \(p,s\in\mathbf Z_{>0}\), then
\[
       \frac{1}{2as}\log|u^p|
          =\frac{p}{as}\log|t|
\]
is the required extension; increasing rational approximation to \(-c\)
gives the general finite negative case by a decreasing sequence.  Thus the
mixed-zero obstruction is genuinely multibranch and is absent on this
unibranch singularity.
\end{example}

The ordinary multinode makes the conductor construction completely
explicit and will be useful as a permanent boundary example.

\begin{theorem}[Radial psh classification on an ordinary multinode]
\label{thm:p122-multinode-radial-psh-classification}
Let \(K\) be a complete nontrivially valued non-Archimedean field, choose
\(\varpi\in K\) with \(0<|\varpi|<1\), and put
\(r=|\varpi|\) and \(a=-\log r\).  For \(q\geq2\), set
\[
 A_q:=K\langle x_1,\ldots,x_q\rangle/
             (x_ix_j:i\ne j),
 \qquad U_q:=\mathcal M(A_q),
\]
let \(D_i=V(x_j:j\ne i)\) be the \(i\)-th closed-disc branch, and let
\(o=V(x_1,\ldots,x_q)\).  Then
\(U_q=\bigcup_iD_i\) and the branches meet only at \(o\).  Given
\(c=(c_1,\ldots,c_q)\in[-\infty,0]^q\), define
\(h_c^\circ:U_q\setminus\{o\}\to[-\infty,0]\) on the \(i\)-th branch by
\[
 h_c^\circ(z):=
 \begin{cases}
  -\infty,&c_i=-\infty,\\[1mm]
  -\dfrac{c_i}{a}\log|x_i(z)|,&c_i> -\infty.
 \end{cases}
 \tag{\(\blacktriangle\)}\label{eq:p122-multinode-radial-function}
\]
At the radius-\(r\) Gauss point of \(D_i\), this function has value \(c_i\).
The following are equivalent:
\begin{enumerate}
\item \(h_c^\circ\) extends to a classically psh function on \(U_q\);
\item \(h_c^\circ\) extends to a CLD-v2 psh function on \(U_q\);
\item either \(c=0\), or \(c_i<0\) for every \(i\), with
      \(-\infty<0\) allowed.
\end{enumerate}
When it exists, the extension is unique.  Its nodal value is \(0\) when
\(c=0\), and \(-\infty\) otherwise.

Consequently, after identifying the branch link with \(\{1,\ldots,q\}\),
the radial-extendable part of the completed ideal-psh orthant is exactly
\[
                 \{0\}\ \cup\ [-\infty,0)^q.
 \tag{\(\blacklozenge\)}\label{eq:p122-multinode-radial-psh-cone}
\]
It is the zero profile together with the pointwise decreasing closure of
the primary principal-ideal profiles.  In particular, the mixed-zero
profile in \Cref{ex:p122-curve-ideal-psh-boundary} is excluded, whereas
every strictly negative finite profile and every limit obtained by sending
some of its coordinates to \(-\infty\) has a radial CLD-v2 psh extension.
This classifies only the displayed radial realization on the ordinary
multinode; it does not identify arbitrary boundary traces or settle P12.2.
\end{theorem}

\begin{proof}
Assume first that \(c=0\).  The constant zero function is a CLD-v2 psh
extension.  Now suppose that every \(c_i<0\), and put
\[
 I:=\{i:c_i> -\infty\},\qquad J:=\{i:c_i=-\infty\}.
\]
If \(I=\varnothing\), the constant \(-\infty\) function is a CLD-v2 psh
extension.  Otherwise, put \(\lambda_i=-c_i\in(0,+\infty)\) for
\(i\in I\).  Choose positive rational numbers
\[
       \lambda_i^{(n)}=k_{n,i}/N_n\nearrow\lambda_i
       \qquad(i\in I),
\]
using one positive denominator \(N_n\) at each stage, and set
\[
                  f_n:=\sum_{i\in I}x_i^{k_{n,i}}\in A_q,
       \qquad g_n:=(aN_n)^{-1}\log|f_n|.
\]
On \(D_i\) for \(i\in I\), all summands except the \(i\)-th vanish, and
hence
\[
                  g_n|_{D_i}=a^{-1}\lambda_i^{(n)}\log|x_i|.
\]
On \(D_j\) for \(j\in J\), the function \(f_n\) vanishes identically, so
\(g_n|_{D_j}=-\infty\).  In particular, omitting a summand---rather than
sending its exponent to infinity---is what gives the prescribed value at
the outer points with \(|x_j|=1\).
Each \(g_n\) is CLD-v2 psh by
\Cref{prop:ch15-log-approachable,prop:ch15-psh-cone}.  Since
\(\log|x_i|\leq0\), the sequence \((g_n)\) decreases pointwise on every
branch.  Its limit is \(h_c^\circ\) off \(o\), takes the value \(-\infty\)
at \(o\), and is CLD-v2 psh by decreasing closedness.  This proves
(3)\(\Rightarrow\)(2), and (2)\(\Rightarrow\)(1) is
\Cref{thm:cld-to-classical-psh}.

For the converse, suppose that a classically psh extension \(h\) exists and
that \(c_i=0\) for some \(i\).  The rigid points of \(D_i\setminus\{o\}\)
with \(x_i=\varpi^m\) converge to \(o\), so upper semicontinuity gives
\(h(o)\geq0\).  If \(c\ne0\), choose \(j\) with \(c_j<0\).  On the closed
radius-\(r\) disc
\[
                    D_{j,r}:=\{z\in D_j:|x_j(z)|\leq r\},
\]
the unique Shilov point is its radius-\(r\) Gauss point, where \(h=c_j<0\).
When \(c_j\) is finite, every other nonnodal value on this disc is at most
\(c_j\); when \(c_j=-\infty\), all its nonnodal values are \(-\infty\).
Thus \(h(o)\geq0\) contradicts the affinoid maximum principle
\cite[Theorem~6.4]{GR25Classical}.  This proves (1)\(\Rightarrow\)(3).

It remains only to prove uniqueness at the node.  For \(c=0\), upper
semicontinuity gives \(h(o)\geq0\), while the same maximum principle on any
branch disc gives \(h(o)\leq0\).  If all \(c_i<0\) and one \(c_j\) is
finite, applying the maximum principle on the radius-\(r^m\) disc in
\(D_j\) gives \(h(o)\leq m c_j\) for every \(m\), hence
\(h(o)=-\infty\).  If every \(c_i=-\infty\), one branch disc gives the same
conclusion immediately.

Finally, strictly negative rational vectors in \(\mathbf Q_{<0}^q\) are
exactly the radius-\(r\) profiles of the functions
\((aN)^{-1}\log|\sum_i x_i^{k_i}|\).  Coordinatewise
increasing rational approximation of
\((\lambda_i)\in(0,+\infty]^q\) therefore shows that
\([-\infty,0)^q\) is precisely their decreasing pointwise closure.
This last closure statement is on the finite radius-\(r\) branch link; for
the analytic extension of an infinite coordinate, the preceding construction
omits the corresponding summand so that the whole branch, including its
outer boundary, has value \(-\infty\).
The zero profile must be adjoined separately, since every primary
principal-ideal profile is strictly negative on every branch.
\end{proof}
